\section{Lessons from a physical-deployment attempt}
\label{app:real-robot}

To scope what separates the simulation results from a deployable
system, we ran one bring-up session on a physical SO-ARM101
(Feetech STS3215 servos, wrist-mounted RealSense D405, Claude
Opus~4.8 as the perception VLM). The same Generate phase
synthesized a real-robot driver from the SO-101 MJCF/URDF and a
one-line spec. The arm reached a kinesthetically taught grasp pose
within 2\,mm after a one-shot servo-bias correction, and the VLM
localized the target cube in the eye-in-hand view; end-to-end
autonomous grasping of a 4\,cm cube still failed. The
failures break down into four independent gaps.

\noindent\textbf{(1) Servo reproducibility.} Commanded vs.\ actual
joint positions differ by 1--4$^\circ$ per joint under gravity
load, propagating to 5--8\,mm of end-effector error. A
feed-forward bias vector recovers 2\,mm accuracy but is
pose-dependent: one compensation vector does not generalize across
the workspace.

\noindent\textbf{(2) Hand-eye calibration.} The wrist-mounted
camera-to-base transform must be calibrated per arm. Our
perturbation-based grid searches were ill-conditioned (13--17\,cm
predicted-depth error); the working industry practice is a
human-placed calibration target with a stored affine solve, which
is automatic only \emph{after} the human places markers.

\noindent\textbf{(3) Depth sensing at close range.} With the
end-effector at 24\,cm, scene depths of 6--15\,cm sit at the lower
edge of the D405's working range; the gripper fingers themselves
read as invalid depth, and glare on metal surfaces corrupts
detection boxes.

\noindent\textbf{(4) Grasp verification.} The physical arm has no
force, tactile, or slip sensing; grasp success must be inferred
from a single gripper-encoder integer. That signal is too coarse
to drive the outer verify-and-retry loop that every simulation
result in this paper relies on.

None of the four gaps is specific to LLM-written code; a
hand-written driver would face the same per-arm bring-up. What
simulation removes is precisely the feedback channel
(\S\ref{sec:verifier}) that Auto-Adapter's validation depends on.
The practical boundary is therefore: simulator-grounded synthesis
automates the \emph{algorithmic} layer of onboarding (kinematics,
control, primitive API), while calibration, bias compensation, and
contact sensing remain a per-robot engineering layer that today
takes a human a session of bring-up time.

\paragraph{Safety requirements for any physical extension.}
\label{app:safety}
In this paper the write--execute--revise loop never touches
hardware: all generation, execution, and revision happen in MuJoCo,
and the exported driver is a static, auditable artifact. The
bring-up session above used conservative manual operation with
z-bounds and joint-limit checks in the bridge. Before generated
code executes on hardware we consider the following minimally
necessary: (i) non-overridable driver-level clamps on joint limits,
velocity, and torque; (ii) a workspace bounding volume enforced
outside the generated code; (iii) sim-first re-validation of any
revised driver before hardware execution; (iv) collision and force
monitoring with automatic stop; (v) a hardware e-stop; and (vi) a
human approval gate with staged, slowed rollout for the first
hardware run of any newly generated driver.

\paragraph{Future work: scaling physical validation.}
Our concrete next step is perception-grounded reach and place on
this physical SO-ARM101 with the wrist-mounted RGB-D camera,
executed under the safety requirements above (sim-first
re-validation, driver-level clamps, human approval gate), taking
the four gaps documented in this appendix as the work items; a
controlled study of an engineer using LLM coding assistants for
the same onboarding task is a second open comparison.
